\documentclass[11pt,a4paper]{article}

\usepackage[utf8]{inputenc}
\usepackage[T1]{fontenc}
\usepackage{lmodern}
\usepackage{amssymb}
\usepackage[margin=2.4cm]{geometry}
\usepackage{graphicx}
\usepackage{booktabs}
\usepackage{longtable}
\usepackage{array}
\usepackage{enumitem}
\usepackage{xcolor}
\usepackage{tcolorbox}
\usepackage{fancyhdr}
\usepackage{listings}
\usepackage{titlesec}
\usepackage[hidelinks]{hyperref}

\definecolor{asamblue}{RGB}{31,90,140}
\definecolor{asamgrey}{RGB}{95,105,115}
\definecolor{codebg}{RGB}{245,246,248}
\definecolor{warnbg}{RGB}{255,246,230}
\definecolor{warnline}{RGB}{204,132,0}
\definecolor{tipbg}{RGB}{237,247,240}
\definecolor{tipline}{RGB}{46,125,80}

\titleformat{\section}{\Large\bfseries\color{asamblue}}{\thesection}{0.7em}{}
\titleformat{\subsection}{\large\bfseries\color{asamblue}}{\thesubsection}{0.7em}{}

\lstset{
  basicstyle=\ttfamily\small,
  backgroundcolor=\color{codebg},
  frame=single,
  rulecolor=\color{asamgrey!40},
  breaklines=true,
  columns=fullflexible,
  xleftmargin=0.4em,
  framexleftmargin=0.4em,
  literate={-}{{-}}1 {~}{{\textasciitilde}}1
}

\newtcolorbox{warnbox}{colback=warnbg,colframe=warnline,boxrule=0.8pt,
  left=6pt,right=6pt,top=5pt,bottom=5pt}
\newtcolorbox{tipbox}{colback=tipbg,colframe=tipline,boxrule=0.8pt,
  left=6pt,right=6pt,top=5pt,bottom=5pt}

\pagestyle{fancy}
\fancyhf{}
\lhead{\small\color{asamgrey}Agentic-SAM v2 --- User Manual}
\rhead{\small\color{asamgrey}\thepage}
\renewcommand{\headrulewidth}{0.3pt}

\newcommand{\ui}[1]{\textbf{#1}}
\newcommand{\code}[1]{\texttt{\small #1}}

\begin{document}

\begin{titlepage}
\centering
\vspace*{3cm}
{\Huge\bfseries\color{asamblue} Agentic-SAM v2\par}
\vspace{0.6cm}
{\LARGE An AI Assistant for Ultrasound-Guided\\[0.2cm] Radiofrequency Ablation of PTMC\par}
\vspace{1.2cm}
{\large User Manual\par}
\vspace{2cm}
{\large\itshape Real-time PTMC segmentation and tracking $\cdot$ ablation-coverage
monitoring $\cdot$ conversational surgical assistance\par}
\vfill
\begin{warnbox}
\textbf{Intended use.} This software is a research and workflow-prototyping tool.
It provides \emph{decision support only}. It does not diagnose, does not measure
thermal ablation, and must never be used as the sole basis for an
intraoperative decision. Every overlay, measurement, and report must be
verified visually by the operating surgeon.
\end{warnbox}
\vspace{1cm}
{\large\today\par}
\end{titlepage}

\tableofcontents
\newpage

%=======================================================================
\section{What This Application Does}
%=======================================================================

Agentic-SAM v2 is an AI assistant that watches the intraoperative ultrasound
stream alongside the surgeon during radiofrequency ablation (RFA) of a
papillary thyroid microcarcinoma (PTMC). The surgeon marks the lesion once;
from that moment the system segments and follows it frame by frame, measures
how the visible target changes, keeps a timestamped record of the procedure,
and answers questions in natural language --- all on local hardware, with no
patient data leaving the machine.

\subsection{Feature Summary}

\begin{longtable}{@{}p{4.6cm}p{10.4cm}@{}}
\toprule
\textbf{Feature} & \textbf{What it does} \\
\midrule
\endhead
Click-to-segment targeting &
The surgeon pauses the video on any frame and clicks the PTMC. MedSAM2
segments it on that frame immediately. \\
\addlinespace
Continuous lesion tracking &
Playback resumes with the segmentation overlaid, updated on every subsequent
frame, so the lesion stays outlined while the probe and anatomy move. \\
\addlinespace
Live AI monitor &
Every few seconds a short situational note appears: tracking stability, size
changes, and what may be worth re-checking visually. \\
\addlinespace
Ablation-coverage monitoring &
The first segmented frame is stored as the baseline lesion. Coverage and
residual percentages are computed continuously as ablation proceeds. \\
\addlinespace
Size-drop alerts &
When the visible target shrinks substantially relative to the session
maximum, an alert is written to the procedure log with an estimated
percentage. \\
\addlinespace
Conversational assistance &
Questions can be asked at any time, about the current frame or the procedure
so far. Tracking pauses while the assistant answers, then resumes. \\
\addlinespace
Text-prompted segmentation &
Structures can also be segmented from a phrase such as \code{thyroid nodule}
or \code{ablation zone}, without drawing a box. \\
\addlinespace
Procedure event log &
Phase changes, tracking events, alerts, snapshots, and questions are recorded
with timestamps and the RFA phase in force. \\
\addlinespace
Ablation report &
One button produces an illustrated report: baseline lesion, final ablation
pattern, coverage over time, residual assessment, limitations. \\
\addlinespace
Session summary note &
A separate button drafts a structured technical summary of the whole session
from the event log. \\
\bottomrule
\end{longtable}

\subsection{How an RFA Surgeon Benefits}

\paragraph{The target stays outlined while you work.}
PTMC lesions are small and low-contrast, and during ablation the picture
degrades: gas and hyperechoic change obscure margins, and probe pressure
shifts the plane. Marking the lesion once and having it outlined continuously
reduces the effort of repeatedly re-identifying it on a changing image.

\paragraph{Coverage becomes visible instead of remembered.}
Judging which parts of a lesion have already been treated is normally done
from memory of where the electrode has been. The system accumulates a
\emph{coverage map} across the whole session and reports the percentage of the
baseline lesion that appears treated and the percentage that does not ---
turning a mental tally into an explicit, reviewable picture.

\paragraph{Residual regions are pointed out, not just counted.}
If part of the baseline lesion never shows the change associated with
treatment, it is highlighted in red on the ablation map and its location is
described in the report --- a prompt to visually re-check that region before
finishing.

\paragraph{An assistant that answers rather than distracts.}
Questions such as ``how much is ablated?'' or ``did I cover the whole
lesion?'' are answered from measured data in one short reply. Answers are
deliberately terse: an operating room has no time for prose.

\paragraph{Documentation writes itself.}
Every notable event is timestamped during the case. The post-procedure report
is generated from that record, so producing documentation costs one button
press rather than recall after the fact.

\paragraph{Everything runs locally.}
Segmentation models and the language model both run on the workstation's own
GPU. No images or clinical text are transmitted off the machine.

\begin{warnbox}
\textbf{How to read ``coverage''.} Coverage is a \emph{tracking-based proxy}.
The software counts a region as treated when the tracker stops capturing it,
because ablated tissue becomes hyperechoic and loses its previous appearance.
That signal correlates with treatment but is not a measurement of thermal
dose, and it can also be produced by probe motion, plane change, or tracking
error. Treat coverage figures as a prompt to look, never as proof of
treatment.
\end{warnbox}

%=======================================================================
\section{System Overview}
%=======================================================================

The application runs as a \emph{single process} that serves the web interface
and holds the segmentation models in memory, plus one local inference server
for the language model.

\begin{longtable}{@{}p{4.2cm}p{10.8cm}@{}}
\toprule
\textbf{Component} & \textbf{Role} \\
\midrule
\endhead
MedSAM2 & Segments the lesion from the surgeon's box and tracks it through
subsequent frames. This is the model behind every overlay on the Live
Procedure page. \\
\addlinespace
Medical-SAM3 & Segments structures from a text prompt on medical ultrasound. \\
\addlinespace
SAM 3.1 & General-purpose text-prompted segmentation, used as a fallback when
Medical-SAM3 returns nothing. \\
\addlinespace
Language model (vLLM) & Runs locally through an OpenAI-compatible interface.
Powers the live monitor notes, the conversational answers, and both reports. \\
\addlinespace
Agent (LangGraph) & Decides which tool to use for a question, remembers the
conversation, and grounds safety-critical answers in measured data. \\
\bottomrule
\end{longtable}

All models are loaded \emph{before} the interface becomes reachable, so the
first segmentation of a case is as fast as every later one --- there is no
cold start in the middle of a procedure.

\subsection{Starting the Application}

\begin{lstlisting}[language=bash]
cd /home/ubuntu/research/agentic-sam-v2
bash run.sh        # uses the settings in the CONFIGURATION block
bash run.sh 2      # ... but runs everything on GPU 2
\end{lstlisting}

Startup prints its progress and takes roughly a minute (longer the first time,
when the language model is downloaded):

\begin{lstlisting}
[serve] LLM ready: vLLM serving Qwen/Qwen2.5-VL-7B-Instruct
[serve] MedSAM2 loaded in 3.8s
[serve] Medical-SAM3 loaded in 17.2s
[serve] SAM 3.1 loaded in 8.5s
[serve] all models warm -- starting UI on 0.0.0.0:7870
\end{lstlisting}

The interface is then available at \code{http://<workstation>:7870}. Press
\ui{Ctrl+C} in the terminal to stop the application and the inference server
together.

\subsection{Configuration}

The top of \code{run.sh} contains a configuration block; edit it rather than
passing environment variables.

\begin{longtable}{@{}p{4.4cm}p{3.4cm}p{7.2cm}@{}}
\toprule
\textbf{Setting} & \textbf{Default} & \textbf{Meaning} \\
\midrule
\endhead
\code{GPU} & \code{0} & GPU the whole stack uses; the first command-line
argument overrides it. \\
\code{VLLM\_MEM\_FRACTION} & \code{0.35} & Share of that GPU reserved for the
language model; the segmentation models use the remainder. \\
\code{APP\_ENV} & agentic env & Python environment running the interface and
segmentation models. \\
\code{VLLM\_ENV} & vllm\_m3 env & Python environment running the inference
server. \\
\code{LLM\_MODEL} & Qwen2.5-VL-7B & Language model served locally. Must
support images and tool calling. \\
\code{UI\_PORT} & \code{7870} & Port the interface listens on. \\
\bottomrule
\end{longtable}

Two further settings are worth knowing, set in the \code{.env} file:

\begin{itemize}[leftmargin=1.4em]
  \item \code{ASAM\_PIXEL\_SPACING\_MM} --- millimetres per pixel for the
  ultrasound preset in use. When set, all measurements are additionally
  reported in mm and mm\textsuperscript{2}; when left at 0, measurements are
  in pixels only.
  \item \code{ASAM\_ABLATION\_SAMPLE\_SECONDS} --- how often the coverage
  monitor samples the tracked lesion (default: every 5 seconds).
\end{itemize}

%=======================================================================
\section{The Live Procedure Page}
%=======================================================================

This is the page used during the case. It shows the ultrasound frame with the
segmentation overlay, the live assistant, the coverage metrics, and the
recent-events list.

\subsection{Step-by-Step Workflow}

\begin{enumerate}[leftmargin=1.6em]
  \item \textbf{Load the video.} Choose \ui{Sample RFA video} or \ui{Upload}
  and select a file. Frames are extracted once; the frame count and rate are
  reported underneath.

  \item \textbf{Set the RFA phase.} The \ui{RFA phase} selector offers
  \emph{Pre-op review}, \emph{Targeting}, \emph{Ablation}, and
  \emph{Post-ablation assessment}. The current phase is stamped on every event
  recorded from then on, and is given to the assistant as context.

  \item \textbf{Find the lesion.} Press \ui{$\blacktriangleright$ Play} to run
  the video, and \ui{$\parallel$ Pause} when the PTMC is well seen. The
  \ui{Sampled frame} slider can also be dragged to any frame; dragging it
  pauses playback automatically.

  \item \textbf{Mark the PTMC.} Click on the lesion in the paused frame. A
  marker and a square target box appear. Adjust \ui{Tracking box size (px)}
  beforehand if the default box does not enclose the lesion, and click again
  to reposition.

  \item \textbf{Start tracking.} Press \ui{Track PTMC}. The lesion is
  segmented on the marked frame and playback resumes automatically, with the
  overlay redrawn on each new frame as it is segmented. The manual marker
  disappears --- the live overlay replaces it.

  \item \textbf{Ablate.} Perform the ablation while the overlay follows the
  lesion. Coverage and residual percentages update as treatment proceeds, and
  the live assistant comments on what it observes.

  \item \textbf{Pause whenever needed.} \ui{$\parallel$ Pause} halts tracking;
  \ui{$\blacktriangleright\!\!\mid$ Continue} resumes it from where it
  stopped. While paused, clicking the image again re-marks the lesion, and
  pressing \ui{Track PTMC} restarts tracking from the corrected position ---
  useful after a large probe movement.

  \item \textbf{Capture anything notable.} \ui{Snapshot to event log} stores
  the current overlaid frame together with the coverage figures.
\end{enumerate}

\begin{tipbox}
\textbf{Choosing where to start tracking.} Start from a frame where the lesion
margin is clearly visible. Tracking quality throughout the case depends
heavily on the quality of that first segmentation --- it also becomes the
baseline against which all coverage is measured.
\end{tipbox}

\subsection{Reading the Right-Hand Panel}

\paragraph{Live assistant.} A short note refreshed every few seconds: a
one-line status followed by at most three observations. It reports tracking
stability and apparent size changes, and suggests visual checks. It never
gives instructions.

\paragraph{Tracked PTMC area.} The current segmented area in pixels, with its
change relative to the largest area seen during the session.

\paragraph{Ablation coverage and residual PTMC.} The percentage of the
baseline lesion that appears treated, and the percentage that does not. Both
are marked with an asterisk as proxies. The caption shows how many samples
have been taken.

\paragraph{Recent events.} The last few entries of the procedure log, each
with its time and type: tracking events, alerts, snapshots, questions, and
phase changes.

\subsection{Asking Questions During the Procedure}

Below the video, \ui{Ask during the procedure} accepts a typed question or one
of three preset questions:

\begin{itemize}[leftmargin=1.4em]
  \item \emph{Is the PTMC tracking stable in this frame?}
  \item \emph{Has the visible PTMC area changed recently?}
  \item \emph{What should be visually re-checked right now?}
\end{itemize}

Asking pauses tracking; the assistant answers using the current frame and the
procedure context, and the exchange is added to the event log. Press
\ui{Continue} to resume. Ablation questions --- coverage, residual, how much
is treated --- are always answered from the measured coverage data rather than
from the model's impression of the image.

%=======================================================================
\section{The Image Analysis Page}
%=======================================================================

This page is for examining a single still image: a pre-operative capture, a
frame of interest, or any exported ultrasound image. It is a conversation
rather than a live feed.

\begin{enumerate}[leftmargin=1.6em]
  \item Choose \ui{Sample PTMC image} or \ui{Upload}.
  \item Optionally click two opposite corners to draw a box around the lesion.
  The box is shown with its coordinates; \ui{Clear box} removes it.
  \item Ask a question by typing it or selecting a preset (\emph{Where is the
  PTMC in this ultrasound image?}, \emph{Describe this ultrasound image.},
  \emph{How large is the segmented target?}).
\end{enumerate}

The assistant chooses its own approach: with a box drawn and a localization
question asked, it segments inside the box with MedSAM2; without a box it
segments from the text of the question; for size questions it measures the
most recent segmentation; for background questions it may consult the web.
The overlay appears in the conversation and in \ui{Latest segmentation
overlay}, and each answer carries an expandable \ui{Tools used} record showing
exactly which tool produced the numbers.

The conversation has memory: follow-up questions such as ``and how large is
it?'' are understood in the context of what was just discussed. The
conversation is shared with the Live Procedure page, so a question asked
during the case and a follow-up asked here belong to the same thread.

%=======================================================================
\section{The Procedure Report Page}
%=======================================================================

This page reviews the completed case and produces documentation. Everything it
shows is reconstructed from the event log, so it can be opened mid-case as
well.

\subsection{Session Overview}

Counters for events, snapshots, the final phase, and the current area relative
to the session maximum; a chart of tracked lesion area over time; the full
event table (time, phase, type, message); and a gallery of all snapshots.

\subsection{Ablation Coverage}

\begin{itemize}[leftmargin=1.4em]
  \item \textbf{Metrics} --- final coverage and residual percentages, the
  baseline lesion area, and the number of samples taken.
  \item \textbf{Baseline PTMC} --- the first tracked frame with its
  segmentation, i.e.\ the lesion as it appeared before treatment.
  \item \textbf{Ablation pattern} --- the coverage map. \textcolor{tipline}
  {\textbf{Green}} marks baseline regions counted as treated;
  \textcolor{red}{\textbf{red}} marks regions never observed to change, which
  may not have been covered.
  \item \textbf{Coverage chart} --- how coverage accumulated over the case,
  revealing whether treatment progressed steadily or plateaued.
\end{itemize}

\subsection{Generating Reports}

\ui{Generate ablation report} produces an illustrated analysis: ablation
extent over time, the spatial pattern, an assessment of residual regions and
where they lie, supporting signals such as echo-brightness change, and the
limitations of the method. \ui{Generate session summary note} produces a
structured technical summary of the whole session from the event log. Both
can be downloaded as Markdown.

\begin{tipbox}
\textbf{Why reports are generated on demand.} All monitoring during the case
is ordinary image arithmetic and costs nothing in language-model terms. The
language model is invoked once per report, when the button is pressed. This
keeps the live workflow fast and inference cost low.
\end{tipbox}

%=======================================================================
\section{The System Status Page}
%=======================================================================

For checking the installation and diagnosing problems. It lists each model
with its checkpoint, size, and whether it is loaded; offers manual preload
buttons; shows GPU utilisation and memory; and reports the session
identifier, the language-model endpoint, and the active settings including
pixel spacing.

If a model shows as absent, its weights are missing --- run
\code{python downloader.py} to restore them. If the language-model status is
not \emph{ready}, the inference server did not start; check
\code{.logs/vllm.log}.

%=======================================================================
\section{Troubleshooting}
%=======================================================================

\begin{longtable}{@{}p{5.2cm}p{9.8cm}@{}}
\toprule
\textbf{Symptom} & \textbf{Cause and remedy} \\
\midrule
\endhead
The video will not load and a Git LFS message appears &
The file is a placeholder stub, not a real video. Copy the actual recording
into \code{data/video/} or upload a different file. \\
\addlinespace
Tracking stops early &
The tracker reached the end of the sampled frames, or lost the target. Check
the status line under the video; re-mark the lesion and press
\ui{Track PTMC} again. \\
\addlinespace
The overlay drifts off the lesion &
Probe or plane movement outstripped the tracker. Pause, click the lesion
again, and restart tracking; coverage accumulated so far is retained. \\
\addlinespace
Coverage rises without ablation &
The tracker is losing the target for other reasons (motion, shadowing).
Coverage is a proxy; re-mark the lesion and interpret the figures with that
in mind. \\
\addlinespace
Live assistant shows an error &
The inference server is unreachable. Check the System Status page and
\code{.logs/vllm.log}. Tracking and segmentation continue to work without
it. \\
\addlinespace
\code{FileNotFoundError: 'ninja'} at startup &
The inference environment's \code{bin} directory is not on \code{PATH}.
\code{run.sh} handles this; ensure \code{VLLM\_ENV} points at the environment
that actually contains vLLM. \\
\addlinespace
Out-of-memory on startup &
Lower \code{VLLM\_MEM\_FRACTION} in \code{run.sh}, or select a GPU with more
free memory. \\
\addlinespace
GPU memory stays occupied after exit &
Stop with \ui{Ctrl+C} rather than closing the terminal, so the inference
server is shut down with the application. \\
\bottomrule
\end{longtable}

%=======================================================================
\section{Safety and Limitations}
%=======================================================================

\begin{enumerate}[leftmargin=1.6em]
  \item \textbf{Decision support only.} The system does not diagnose and does
  not recommend treatment. The surgeon remains responsible for every decision.

  \item \textbf{Coverage is a proxy.} It is inferred from loss of tracking
  against the baseline lesion, not from thermal measurement. Probe motion,
  plane change, shadowing, and tracking error can all produce the same signal.

  \item \textbf{Measurements are pixel-based} unless a pixel spacing has been
  configured for the ultrasound preset in use. Millimetre values are only as
  accurate as that calibration.

  \item \textbf{Segmentation can fail silently.} A confidently drawn outline
  may still be wrong. Visual verification is required at every stage.

  \item \textbf{Reports are technical summaries}, generated from what the
  software recorded. They are not clinical operative notes and are not a
  substitute for them.

  \item \textbf{Only what is visible is assessed.} Lesion extent outside the
  imaging plane, and any tissue not in view, are invisible to the system and
  are not represented in coverage figures.
\end{enumerate}

\vspace{1em}
\begin{warnbox}
Not a medical device. Not cleared or approved by any regulatory body. For
research and workflow prototyping only.
\end{warnbox}

\end{document}
